{\bf Back of the envelope}\\
In addition to solving the problems below, discuss the typical values and/or
ranges for the parameters you will need?
\begin{enumerate}
    \item How long does it take a thermal perturbation at the base of a
    craton to reach the surface?  Consider what you will need to know to
    solve this problem.  
    
    {\bf Answer:} \\
    Again, we will use the thermal length estimate for the time
    (Equation~\ref{eq:tscale}).  We need only two parameters then to
    estimate how long it takes the perturbation to reach the surface: the
    thermal diffusivity and the thickness of the lithosphere.  The
    lithosphere is typically between 70 and 200~km thick.  The thermal
    diffusivity of most Earth materials are about $1 \times
    10^{-6}$~m$^{2}$~Ma$^{-1}$ which we will want in units of
    km$^{2}$~Ma$^{-1}$ so that we arrive at the correct units of Ma in our
    final result.
    
    So a little trick here.  If you want to convert the number of seconds
    into years, you can approximate the number of seconds in a year as $\pi
    \times 10^7$.  The real number is 31557600 so our approximation is in
    error by 0.45\% (who cares).  This can be a really handy trick if you
    want to cancel a $\pi$ out of a back the envelope equation.

    Now back to the problem, converting the thermal diffusivity,
    \begin{align*}
        \kappa &= (10^{-6} \text{ m$^{2}$ s$^{-1}$})(\pi \times
        10^{13}
        \text{ s Ma$^{-1}$})(10^{-6} \text{ km m$^{-1}$}) \\
        &\approx 30 \text{ km$^2$ Ma$^{-1}$} .
    \end{align*}
    Inputting our diffusivity and lithospheric thickness into the thermal
    length equation,
    \begin{align*}
        t &= \frac{(70\text{ km})^2}{(4) (30 \text{ km$^2$~Ma$^{-1}$})}\\
        &\approx 40 \text{ Ma},
    \end{align*}
    if the lithosphere is thin and 
    \begin{align*}
        t &= \frac{(200\text{ km})^2}{(4) (30 \text{ km$^2$~Ma$^{-1}$})}\\
        &\approx 330 \text{ Ma},
    \end{align*}
    if the lithosphere is thick.  Either way, it can take a very long time
    for a temperature perturbation to reach the surface.

    \item The retreat of glaciers at the end of the Pleistocene
    ($\sim$10~ka) allowed warmer air temperatures to heat the ground.  How
    deep a borehole would we need to observe temperatures undisturbed by
    warming? Assume a diffusivity of $0.8 \times 10^{-6}$~m$^2$~s$^{-1}$.

    {\bf Answer:} \\
    This time, we know the time so we and the diffusivity so all we need now
    is to estimate the depth the heat pulse has penetrated, which we
    estimate using Equation~\ref{eq:zscale}
    \begin{align*}
        z &= 2 [(0.8 \times 10^{-6} \text{ m$^{2}$~s$^{-1}$})(\pi \times
        10^{11} \text{ s ka$^{-1}$})(10 \text{ ka})]^{1/2} \\
        z &\approx 1000 \text{ m}
    \end{align*}
    So in 10,000 years, the change in surface temperature that resulted from
    glacial retreat has propagated only about 1000~m (1~km) into the Earth.
\end{enumerate}

